\section{Verifica Geometrica della Calibrazione}
\label{sec:verifica_geometrica}

Al termine della calibrazione, oltre a valutare
\(\epsilon_{\mathrm{RMSE}}\), è possibile verificare visivamente la
qualità dei parametri stimati in due modi complementari: la
\emph{proiezione del sistema di riferimento} e la \emph{rettifica
prospettica} dell'immagine.

%─────────────────────────────────────────────
\subsection{Proiezione del sistema di riferimento}
\label{sub:proiezione_assi}

Si definisce l'origine \(O=(0,0,0)^T\) e tre punti
\(P_X=(\ell,0,0)^T\), \(P_Y=(0,\ell,0)^T\), \(P_Z=(0,0,\ell)^T\)
nel sistema affine del pattern, con \(\ell\) lunghezza fisica nota.
Tramite la funzione di proiezione
\[
  \lambda\,\homo{x} = \mat{K}\,[\mat{R}_j \mid \mathbf{t}_j]\,\homo{X}
\]
i quattro punti vengono proiettati nel piano immagine e si tracciano
i segmenti \(\overrightarrow{O'P_X'}\), \(\overrightarrow{O'P_Y'}\),
\(\overrightarrow{O'P_Z'}\) con colori RGB\@. Se la calibrazione è
corretta, gli assi convergono coerentemente all'origine del target
e le loro lunghezze apparenti rispettano le proporzioni metriche
del pattern; l'asse \(Z\) (normale al piano) si accorcia
proporzionalmente all'angolo di inclinazione, comportamento atteso
della proiezione prospettica che non conserva le lunghezze.

%─────────────────────────────────────────────
\subsection{Rettifica prospettica e geometria affine}
\label{sub:rettifica}

Il secondo strumento di verifica è la \emph{rettifica}: applicare
la proiettività inversa \(\mat{H}^{-1}\) all'immagine per riportare
il pattern in una vista frontale canonica. Poiché \(\mat{H}\) è una
proiettività invertibile di \(\mathbb{P}^2\), anche \(\mat{H}^{-1}\)
lo è: l'operazione non cambia classe di trasformazione. La rettifica
è corretta quando le rette parallele del pattern --- che nella
fotografia prospettica convergono verso punti di fuga (proiezioni
di punti impropri del piano del pattern) --- tornano effettivamente
parallele nell'immagine raddrizzata, segnale che la struttura affine
del piano è stata recuperata.

Le caratteristiche \emph{zone nere} che compaiono ai bordi
dell'immagine rettificata hanno due cause geometriche: da un lato,
la proiettività \(\mat{H}^{-1}\) può mandare pixel del canvas in
posizioni esterne al frame originale (punti senza colore valido);
dall'altro, la correzione della distorsione radiale e tangenziale
via \((k_1,k_2,k_3,p_1,p_2)\) dilata localmente la griglia di pixel
vicino ai bordi, lasciando regioni non coperte.